% preamble-exam-zh.tex — 共享 preamble
% 用于所有 exam-zh 模板
%
% 样式策略：
%   屏幕 → 语义彩色（蓝=关键想法，红=易错，绿=路线，灰=提示/教师备注）
%   打印 → 自动降级为黑白（通过 print mode 或 monochrome 选项）

\documentclass{exam-zh}

% ---------- 基础配置 ----------
\examsetup{
  page/size = a4paper,
  page/show-head = false,
  page/show-foot = false,
  sealline/show = false,
  font = times,
  math-font = xits,
}

\usepackage{tcolorbox}
\usepackage{needspace}
\usepackage{graphicx}
\usepackage{tikz}
\usetikzlibrary{calc,intersections,angles,quotes,arrows.meta,decorations.markings,patterns.meta,backgrounds}
\usepackage{pgfplots}
\pgfplotsset{compat=1.18}
\tcbuselibrary{skins,breakable}

% ---------- 打印/屏幕颜色切换 ----------
% 用法：\documentclass[monochrome]{exam-zh} 即可切换为纯黑白
% 注意：不要使用 \if@xxx 放在 makeatother 外部；这里改成普通条件名。
\newif\ifeduprintcolor
\eduprintcolortrue

\makeatletter
\@ifclasswith{exam-zh}{monochrome}{\eduprintcolorfalse}{}
\makeatother

% ---------- 语义颜色定义 ----------
% 屏幕：有色彩区分；打印：降级为灰度
\ifeduprintcolor
  % --- 屏幕彩色模式 ---
  \definecolor{edu-blue}{RGB}{47, 82, 122}       % 关键想法
  \definecolor{edu-blue-bg}{RGB}{243, 246, 252}
  \definecolor{edu-red}{RGB}{180, 40, 40}         % 易错提醒
  \definecolor{edu-red-bg}{RGB}{255, 245, 240}
  \definecolor{edu-green}{RGB}{34, 100, 60}       % 路线图
  \definecolor{edu-green-bg}{RGB}{242, 250, 245}
  \definecolor{edu-orange}{RGB}{160, 90, 0}       % 训练目标
  \definecolor{edu-orange-bg}{RGB}{255, 248, 235}
  \definecolor{edu-gray}{RGB}{100, 100, 100}      % 提示/教师备注
  \definecolor{edu-gray-bg}{RGB}{248, 248, 248}
  \definecolor{edu-purple}{RGB}{90, 50, 120}      % 思考题
  \definecolor{edu-purple-bg}{RGB}{248, 244, 252}
\else
  % --- 打印黑白模式 ---
  \definecolor{edu-blue}{gray}{0.15}
  \definecolor{edu-blue-bg}{gray}{0.96}
  \definecolor{edu-red}{gray}{0.25}
  \definecolor{edu-red-bg}{gray}{0.97}
  \definecolor{edu-green}{gray}{0.20}
  \definecolor{edu-green-bg}{gray}{0.97}
  \definecolor{edu-orange}{gray}{0.25}
  \definecolor{edu-orange-bg}{gray}{0.98}
  \definecolor{edu-gray}{gray}{0.40}
  \definecolor{edu-gray-bg}{gray}{0.97}
  \definecolor{edu-purple}{gray}{0.30}
  \definecolor{edu-purple-bg}{gray}{0.98}
\fi
\colorlet{edu-diagram-muted}{edu-gray}

% ---------- TikZ diagram semantic macros ----------
\definecolor{edu-diagram-ink}{RGB}{17,24,39}
\definecolor{edu-diagram-marker}{RGB}{220,38,38}
\definecolor{edu-diagram-angle}{RGB}{5,150,105}
\tikzset{
  diagram point/.style={fill=edu-diagram-ink},
  diagram tick/.style={draw=edu-diagram-marker,line width=0.9pt,line cap=round},
  diagram segment/.style={draw=edu-diagram-ink,line cap=round},
  point label/.style={inner sep=1pt},
  condition label/.style={inner sep=1pt}
}
\providecommand{\DrawSegment}[3][]{\draw[diagram segment,#1] (#2) -- (#3);}
\providecommand{\DrawDashedSegment}[3][]{\draw[diagram segment,dashed,#1] (#2) -- (#3);}
\providecommand{\Triangle}[4][]{\path[#1] (#2) -- (#3) -- (#4) -- cycle;}
\providecommand{\Quadrilateral}[5][]{\path[#1] (#2) -- (#3) -- (#4) -- (#5) -- cycle;}
\providecommand{\PolygonPath}[2][]{\path[#1] #2 -- cycle;}
\providecommand{\DiagramPointRadius}{0.052cm}
\providecommand{\PointDot}[2][]{\fill[diagram point,#1] (#2) circle[radius=\DiagramPointRadius];}
\providecommand{\PointLabel}[3][]{\node[point label,#1] at (#2) {#3};}
\providecommand{\SegmentLabel}[4][]{\node[condition label,#1] at ($(#2)!0.5!(#3)$) {#4};}
\providecommand{\CoordinateTag}[3][]{\node[condition label,#1] at (#2) {#3};}
\providecommand{\AngleMark}[4][]{\pic[draw=edu-diagram-angle,line width=0.9pt,angle radius=0.48cm,#1] {angle=#2--#3--#4};}
\providecommand{\AngleLabel}[5][]{\pic["{#5}",draw=none,angle radius=0.64cm,#1] {angle=#2--#3--#4};}
\providecommand{\RightAngleMark}[4][]{\pic[draw=edu-diagram-marker,line width=0.9pt,angle radius=0.28cm,#1] {right angle=#2--#3--#4};}
\providecommand{\EqualTick}[3][]{%
  \path[postaction={decorate},decoration={markings,mark=at position .5 with {\draw[diagram tick,#1] (0pt,-4pt) -- (0pt,4pt);}}] (#2) -- (#3);%
}
\providecommand{\DoubleEqualTick}[3][]{%
  \path[postaction={decorate},decoration={markings,mark=at position .46 with {\draw[diagram tick,#1] (0pt,-4pt) -- (0pt,4pt);},mark=at position .54 with {\draw[diagram tick,#1] (0pt,-4pt) -- (0pt,4pt);}}] (#2) -- (#3);%
}
\providecommand{\TripleEqualTick}[3][]{%
  \path[postaction={decorate},decoration={markings,mark=at position .43 with {\draw[diagram tick,#1] (0pt,-4pt) -- (0pt,4pt);},mark=at position .5 with {\draw[diagram tick,#1] (0pt,-4pt) -- (0pt,4pt);},mark=at position .57 with {\draw[diagram tick,#1] (0pt,-4pt) -- (0pt,4pt);}}] (#2) -- (#3);%
}
\providecommand{\ParallelMark}[3][]{\EqualTick[#1]{#2}{#3}}
\providecommand{\NamedSegmentPath}[3]{\path[name path=#1] (#2) -- (#3);}
\providecommand{\IntersectPaths}[3]{\path[name intersections={of=#2 and #3, by=#1}];}

% ---------- 自定义教学环境 ----------
% 共性：enhanced + breakable（跨页不断裂）

% 原题卡片 — 强边框，突出显示
\newtcolorbox{problemcard}{
  enhanced, breakable,
  colback=edu-blue-bg,
  colframe=edu-blue,
  boxrule=1.2pt,
  arc=0pt,
  left=3mm, right=3mm, top=3mm, bottom=3mm,
}

% 解题路线图 — 左边线 + 浅底
\newtcolorbox{routebox}{
  enhanced, breakable,
  colback=edu-green-bg,
  colframe=edu-green!30,
  boxrule=0pt,
  borderline west={2.5pt}{0pt}{edu-green},
  left=3mm, right=3mm, top=3mm, bottom=3mm,
}

% 关键想法 — 蓝色左边线
\newtcolorbox{keyideabox}{
  enhanced, breakable,
  colback=edu-blue-bg,
  colframe=edu-blue!30,
  boxrule=0pt,
  borderline west={2.5pt}{0pt}{edu-blue},
  left=3mm, right=3mm, top=3mm, bottom=3mm,
}

% 易错提醒 — 红色左边线
\newtcolorbox{mistakebox}{
  enhanced, breakable,
  colback=edu-red-bg,
  colframe=edu-red!20,
  boxrule=0pt,
  borderline west={3pt}{0pt}{edu-red},
  left=3mm, right=3mm, top=2.5mm, bottom=2.5mm,
}

% 提示 — 灰色虚线左边线
\newtcolorbox{hintbox}{
  enhanced, breakable,
  colback=edu-gray-bg,
  colframe=edu-gray!20,
  boxrule=0pt,
  borderline west={2pt}{0pt}{edu-gray, dashed},
  left=3mm, right=3mm, top=2mm, bottom=2mm,
}

% 思考题 — 紫色虚线左边线
\newtcolorbox{thinkbox}{
  enhanced, breakable,
  colback=edu-purple-bg,
  colframe=edu-purple!20,
  boxrule=0pt,
  borderline west={2pt}{0pt}{edu-purple, dashed},
  left=2.5mm, right=2.5mm, top=2.5mm, bottom=2.5mm,
}

% 教师备注 — 灰色左边线，小字
\newtcolorbox{teachernote}{
  enhanced, breakable,
  colback=edu-gray-bg,
  colframe=edu-gray!15,
  boxrule=0pt,
  borderline west={2pt}{0pt}{edu-gray!60},
  left=2.5mm, right=2.5mm, top=2.5mm, bottom=2.5mm,
  fontupper=\small,
  colupper=black!60,
}

% 训练目标 — 橙色左边线，小字
\newtcolorbox{traininggoal}{
  enhanced, breakable,
  colback=edu-orange-bg,
  colframe=edu-orange!15,
  boxrule=0pt,
  borderline west={1.5pt}{0pt}{edu-orange},
  left=2mm, right=2mm, top=1mm, bottom=1mm,
  fontupper=\small,
}

% 答题区 — 无边框，纯留白
\newcommand{\answerarea}[1][40mm]{%
  \vspace{2mm}%
  \begin{tcolorbox}[
    enhanced, breakable,
    colback=white,
    colframe=white,
    boxrule=0pt,
    height=#1,
    valign=top,
    bottom=0mm,
    after skip=0mm,
  ]
  \end{tcolorbox}%
}

\newsavebox{\diagramtikzbox}

\newenvironment{diagramcoltikz}[2]{%
  \begin{minipage}[t]{#1}
    \vspace{0pt}\centering
    \def\diagramcaption{#2}%
    \begin{lrbox}{\diagramtikzbox}%
}{%
    \end{lrbox}%
    \usebox{\diagramtikzbox}%
    \ifx\diagramcaption\empty\else
      \par\vspace{1mm}{\scriptsize\color{edu-diagram-muted}\diagramcaption}
    \fi
  \end{minipage}%
}

\newenvironment{diagramrowtikz}[3]{%
  \begin{minipage}[t]{#1}
    \vspace{0pt}\centering
    \def\diagramlabel{#2}%
    \def\diagramcaption{#3}%
    \begin{lrbox}{\diagramtikzbox}%
}{%
    \end{lrbox}%
    \usebox{\diagramtikzbox}%
    \ifx\diagramlabel\empty\else
      \par\vspace{1mm}{\scriptsize\bfseries\diagramlabel}
    \fi
    \ifx\diagramcaption\empty\else
      \par\vspace{0.5mm}{\scriptsize\color{edu-diagram-muted}\diagramcaption}
    \fi
  \end{minipage}%
}

\newenvironment{answerareawithdiagramtikz}[3]{%
  \vspace{2mm}%
  \begin{tcolorbox}[
    enhanced, breakable,
    colback=white,
    colframe=white,
    boxrule=0pt,
    height=#1,
    valign=top,
    bottom=0mm,
    after skip=0mm,
  ]
  \begin{minipage}[t]{\dimexpr\linewidth-#2-6mm\relax}
    \vspace{0pt}
  \end{minipage}\hfill
  \begin{diagramcoltikz}{#2}{#3}
}{%
  \end{diagramcoltikz}
  \end{tcolorbox}%
}

% ---------- 字体设置 ----------
% exam-zh (ctexbook) already sets CJK fonts via ctex.
% Only override if specific fonts are needed.
% macOS: Songti SC, STHeiti, STFangsong
% Linux: SimSun, SimHei, FangSong

% ---------- 页面设置 ----------
% exam-zh (ctexbook) already loads geometry.
% Override margins if needed:
% \geometry{top=16mm, bottom=16mm, left=14mm, right=14mm}

\examsetup{
  question/show-answer = true,
  fillin/show-answer = true,
  solution/show-solution = show-stay,
}

% ---------- 教师版解答样式 ----------
\newtcolorbox{solutionblock}{
  enhanced, breakable,
  colback=edu-blue-bg,
  colframe=edu-blue!25,
  boxrule=0pt,
  borderline west={2pt}{0pt}{edu-blue!50},
  arc=0pt,
  left=3mm, right=3mm, top=2mm, bottom=2mm,
  before skip=2mm,
  after skip=3mm,
  fontupper=\small,
}

% 解答步骤编号
\newcounter{solstep}
\newcommand{\solstep}[2]{%
  \stepcounter{solstep}%
  \par\medskip
  \noindent\hangindent=6mm
  {\color{edu-blue}\bfseries\textcircled{\scriptsize\thesolstep}\enspace #1}\enspace #2\par
}

\begin{document}

\section*{矩形存在性综合作业}

\section*{一、固定顺序矩形：从直角入口到轨迹包装（每题 10 分，共 30 分）}

\needspace{8\baselineskip}

\begin{problem}[points=10]
  已知 $A(1,1)$，$B(5,1)$，点 $C$ 在直线 $y=x+1$ 上。若四边形 $ABCD$ 是矩形，求点 $C,D$ 的坐标。

  \answerarea[35mm]
  \setcounter{solstep}{0}
  \begin{solutionblock}
    \textbf{答案：}$C(5,6),D(1,6)$\par\medskip
    \solstep{确定直角}{$AB$ 水平，固定顺序矩形中 $BC\perp AB$，所以 $x_C=5$。}
    \solstep{代入轨迹}{代入 $y=x+1$，得 $C(5,6)$。}
    \solstep{求第四点}{$D=A+C-B=(1+5-5,1+6-1)=(1,6)$。}
  \end{solutionblock}
\end{problem}


\needspace{8\baselineskip}

\begin{problem}[points=10]
  已知 $A(1,2)$，$B(5,3)$，$D(2,m)$。若四边形 $ABCD$ 是矩形，求 $m$ 的值和点 $C$ 的坐标。

  \answerarea[42mm]
  \setcounter{solstep}{0}
  \begin{solutionblock}
    \textbf{答案：}$m=-2$，$C(6,-1)$\par\medskip
    \solstep{写相邻边向量}{$\overrightarrow{AB}=(4,1)$，$\overrightarrow{AD}=(1,m-2)$。}
    \solstep{列垂直关系}{$(4,1)\cdot(1,m-2)=0$，得 $4+m-2=0$，所以 $m=-2$。}
    \solstep{求C}{$C=B+D-A=(5+2-1,3-2-2)=(6,-1)$。}
  \end{solutionblock}
\end{problem}


\needspace{8\baselineskip}

\begin{problem}[points=10]
  已知 $A(0,0)$，$B(4,0)$，点 $C$ 在反比例函数 $y=\dfrac{12}{x}$ 的图像上，且 $x_C>0$。若四边形 $ABCD$ 是矩形，求点 $C,D$ 的坐标。

  \answerarea[40mm]
  \setcounter{solstep}{0}
  \begin{solutionblock}
    \textbf{答案：}$C(4,3),D(0,3)$\par\medskip
    \solstep{确定C}{$AB$ 水平，固定顺序下 $BC$ 竖直，所以 $x_C=4$。代入 $y=\dfrac{12}{x}$ 得 $C(4,3)$。}
    \solstep{求D}{$D=A+C-B=(0,3)$。}
  \end{solutionblock}
\end{problem}


\section*{二、无固定顺序与对角线包装（每题 10 分，共 30 分）}

\needspace{8\baselineskip}

\begin{problem}[points=10]
  已知 $A(1,1)$，$B(5,1)$，点 $C$ 在直线 $y=x+1$ 上。若 $A,B,C,D$ 构成矩形，求所有可能的点 $C,D$。

  \answerarea[55mm]
  \setcounter{solstep}{0}
  \begin{solutionblock}
    \textbf{答案：}$C(1,2),D(5,2)$；或 $C(5,6),D(1,6)$\par\medskip
    \solstep{设点}{设 $C(t,t+1)$。}
    \solstep{直角在A}{$AB$ 水平，$AC\perp AB$ 时 $t=1$，得 $C(1,2)$，$D=B+C-A=(5,2)$。}
    \solstep{直角在B}{$BC\perp AB$ 时 $t=5$，得 $C(5,6)$，$D=A+C-B=(1,6)$。}
    \solstep{直角在C}{列 $(A-C)\cdot(B-C)=0$，得 $2t^2-6t+5=0$，判别式小于 $0$，无实数解。}
  \end{solutionblock}
\end{problem}


\needspace{8\baselineskip}

\begin{problem}[points=10]
  已知 $A(0,0)$，$C(8,6)$，点 $B$ 在直线 $y=3$ 上。若 $A,B,C,D$ 构成以 $AC$ 为对角线的矩形，求所有可能的点 $B,D$。

  \answerarea[55mm]
  \setcounter{solstep}{0}
  \begin{solutionblock}
    \textbf{答案：}$B(9,3),D(-1,3)$；或 $B(-1,3),D(9,3)$\par\medskip
    \solstep{设点与第四点}{设 $B(t,3)$。因 $AC$ 为对角线，$D=A+C-B=(8-t,3)$。}
    \solstep{列直角}{矩形中 $AB\perp BC$，所以 $(t,3)\cdot(8-t,3)=0$。}
    \solstep{求解}{化简得 $t^2-8t-9=0$，解得 $t=9$ 或 $t=-1$。}
    \solstep{写答案}{当 $t=9$ 时，$B(9,3),D(-1,3)$；当 $t=-1$ 时，$B(-1,3),D(9,3)$。}
  \end{solutionblock}
\end{problem}


\needspace{8\baselineskip}

\begin{problem}[points=10]
  已知 $A(0,0)$，$B(6,0)$，点 $C$ 在直线 $y=2x$ 上。若 $A,B,C,D$ 构成矩形，求所有可能的点 $C,D$。

  \answerarea[60mm]
  \setcounter{solstep}{0}
  \begin{solutionblock}
    \textbf{答案：}$C(6,12),D(0,12)$；或 $C\left(\dfrac65,\dfrac{12}{5}\right),D\left(\dfrac{24}{5},-\dfrac{12}{5}\right)$\par\medskip
    \solstep{设点}{设 $C(t,2t)$。}
    \solstep{直角在A}{$AC\perp AB$ 时 $t=0$，此时 $C=A$，舍去。}
    \solstep{直角在B}{$BC\perp AB$ 时 $t=6$，得 $C(6,12)$，$D=A+C-B=(0,12)$。}
    \solstep{直角在C}{列 $(-t,-2t)\cdot(6-t,-2t)=0$，得 $t(5t-6)=0$。舍去 $t=0$，取 $t=\dfrac65$，所以 $C\left(\dfrac65,\dfrac{12}{5}\right)$，$D=A+B-C=\left(\dfrac{24}{5},-\dfrac{12}{5}\right)$。}
  \end{solutionblock}
\end{problem}


\section*{三、参数反求与综合筛选（每题 10 分，共 40 分）}

\needspace{8\baselineskip}

\begin{problem}[points=10]
  已知 $A(0,0)$，$B(6,0)$，点 $C$ 在直线 $y=kx$ 上。若四边形 $ABCD$ 是矩形，且矩形面积为 $36$，求 $k$ 的值以及点 $C,D$ 的坐标。

  \answerarea[45mm]
  \setcounter{solstep}{0}
  \begin{solutionblock}
    \textbf{答案：}$k=1$ 时 $C(6,6),D(0,6)$；$k=-1$ 时 $C(6,-6),D(0,-6)$\par\medskip
    \solstep{确定C的形式}{固定顺序下 $BC\perp AB$，所以 $x_C=6$。又 $C$ 在 $y=kx$ 上，故 $C(6,6k)$。}
    \solstep{用面积}{矩形面积为 $6\cdot |6k|=36$，所以 $|k|=1$。}
    \solstep{写答案}{$k=1$ 时 $C(6,6),D(0,6)$；$k=-1$ 时 $C(6,-6),D(0,-6)$。}
  \end{solutionblock}
\end{problem}


\needspace{8\baselineskip}

\begin{problem}[points=10]
  已知 $A(0,0)$，$B(4,0)$，点 $C$ 在二次函数 $y=x^2+ax+4$ 的图像上。若四边形 $ABCD$ 是矩形，且矩形面积为 $16$，求 $a$ 的值以及点 $C,D$ 的坐标。

  \answerarea[50mm]
  \setcounter{solstep}{0}
  \begin{solutionblock}
    \textbf{答案：}$a=-4$ 时 $C(4,4),D(0,4)$；$a=-6$ 时 $C(4,-4),D(0,-4)$\par\medskip
    \solstep{确定C的横坐标}{$AB$ 水平且固定顺序，所以 $x_C=4$。}
    \solstep{用面积确定纵坐标}{矩形面积为 $4|y_C|=16$，所以 $y_C=4$ 或 $y_C=-4$。}
    \solstep{代入二次函数}{当 $x=4$ 时，$y_C=16+4a+4=20+4a$。若 $20+4a=4$，得 $a=-4$；若 $20+4a=-4$，得 $a=-6$。}
    \solstep{求D}{$a=-4$ 时 $C(4,4),D(0,4)$；$a=-6$ 时 $C(4,-4),D(0,-4)$。}
  \end{solutionblock}
\end{problem}


\needspace{8\baselineskip}

\begin{problem}[points=10]
  已知 $A(0,0)$，$C(6,8)$，点 $B$ 在直线 $y=4$ 上。若 $A,B,C,D$ 构成以 $AC$ 为对角线的矩形，求所有可能的点 $B,D$。

  \answerarea[55mm]
  \setcounter{solstep}{0}
  \begin{solutionblock}
    \textbf{答案：}$B(8,4),D(-2,4)$；或 $B(-2,4),D(8,4)$\par\medskip
    \solstep{设点与第四点}{设 $B(t,4)$。因 $AC$ 为对角线，$D=A+C-B=(6-t,4)$。}
    \solstep{列直角}{$AB\perp BC$，所以 $(t,4)\cdot(6-t,4)=0$。}
    \solstep{求解}{化简得 $t^2-6t-16=0$，解得 $t=8$ 或 $t=-2$。}
    \solstep{写答案}{当 $t=8$ 时，$B(8,4),D(-2,4)$；当 $t=-2$ 时，$B(-2,4),D(8,4)$。}
  \end{solutionblock}
\end{problem}


\needspace{8\baselineskip}

\begin{problem}[points=10]
  已知 $A(0,0)$，$B(4,2)$，点 $C$ 在直线 $y=-2x+k$ 上。若四边形 $ABCD$ 是矩形，且矩形面积为 $20$，求 $k$ 的值以及所有可能的点 $C,D$。

  \answerarea[65mm]
  \setcounter{solstep}{0}
  \begin{solutionblock}
    \textbf{答案：}$k=10$；$C(6,-2),D(2,-4)$，或 $C(2,6),D(-2,4)$\par\medskip
    \solstep{设C并列垂直}{设 $C(x,-2x+k)$。$\overrightarrow{AB}=(4,2)$，$\overrightarrow{BC}=(x-4,-2x+k-2)$。}
    \solstep{确定k}{由 $AB\perp BC$，得 $4(x-4)+2(-2x+k-2)=0$，化简得 $2k-20=0$，所以 $k=10$。}
    \solstep{利用面积}{此时 $C$ 在过 $B$ 且垂直于 $AB$ 的直线 $y=-2x+10$ 上。设 $C=B+s(1,-2)$，则 $BC=|s|\sqrt5$，$AB=2\sqrt5$。}
    \solstep{求C和D}{面积 $20=AB\cdot BC=10|s|$，所以 $s=2$ 或 $s=-2$。当 $s=2$ 时 $C(6,-2),D=A+C-B=(2,-4)$；当 $s=-2$ 时 $C(2,6),D=(-2,4)$。}
  \end{solutionblock}
\end{problem}


\clearpage
\section*{答案速查}




\end{document}
